\section{Conclusão}

Este relatório percorreu conceitos centrais da teoria da probabilidade combinando o
tratamento teórico com experimentos computacionais em R. Em todos os estudos as
estimativas empíricas se aproximaram dos valores previstos pela teoria, e esse
acordo entre o que foi calculado e o que foi observado ilustra como a simulação pode
validar resultados analíticos.

A simulação de experimentos aleatórios mostrou que a frequência relativa de um
evento se estabiliza em torno de sua probabilidade teórica à medida que o número de
repetições cresce. O lançamento da moeda, a retirada de bolas da urna e a soma de
dois dados reproduziram as probabilidades teóricas, e a média amostral da soma
convergiu para o seu valor esperado, como prevê a Lei dos Grandes Números
\cite{ross2010probabilidade}.

O estudo das distribuições discretas e contínuas evidenciou o papel dos principais
modelos na descrição de variáveis aleatórias. A análise das distribuições contínuas
destacou o Teorema Central do Limite, segundo o qual a média de variáveis
independentes e identicamente distribuídas se aproxima de uma distribuição normal
mesmo quando a variável de origem é fortemente assimétrica. A verificação da
esperança e da variância de uma variável exponencial confirmou novamente a Lei dos
Grandes Números, com erro relativo reduzido entre as estimativas e os valores
teóricos.

A probabilidade condicional, a independência e o Teorema de Bayes forneceram a base
para problemas de decisão sob incerteza em aplicações como a filtragem de spam, o
diagnóstico médico, a autenticação biométrica e a detecção de intrusão. Esses
estudos mostraram, em particular, que a probabilidade a posteriori de um evento
depende fortemente de sua prevalência, o que explica por que testes precisos podem
produzir muitos alarmes falsos quando o evento é raro.

A atividade integradora final reuniu esses conceitos em diferentes estudos de
inspiração computacional. O filtro bayesiano de spam, construído sobre a hipótese de
independência condicional, alcançou alta acurácia na separação entre mensagens
legítimas e indesejadas e permitiu ajustar o limiar de decisão conforme a tolerância
a cada tipo de erro. A análise de confiabilidade apontou que a redundância em
paralelo eleva o tempo médio até a falha de um sistema, ao passo que o arranjo em
série o reduz, o que evidencia o ganho obtido com a replicação de componentes. A
modelagem da fila de atendimento revelou que a variabilidade dos tempos de chegada e
de serviço impede que um servidor opere na sua capacidade nominal máxima, gerando
congestionamentos mesmo quando a taxa média de atendimento supera a de chegada.
Estudos como os de confiabilidade e de filas apoiam-se em modelos clássicos da
probabilidade aplicada \cite{ross2024models}.

De modo geral, os resultados indicam que a simulação computacional é um recurso útil
para compreender e validar a teoria da probabilidade, sobretudo em situações de
difícil tratamento analítico \cite{harcholbalter2024probability}. A recorrência da
Lei dos Grandes Números e do Teorema Central do Limite ao longo dos estudos reforça
o papel desses dois resultados, que sustentam tanto a estimação de probabilidades
por frequência quanto o comportamento das médias amostrais observado em cada
experimento.
